% Arvore do repositorio (estado de de9bc41). Pacote dirtree.
\renewcommand*\DTstylecomment{\small\rmfamily\itshape\color{black!65}}
\renewcommand*\DTstyle{\small\ttfamily}
\dirtree{%
.1 morphe/.
.2 README.md\DTcomment{ponto de entrada}.
.2 morphe-up.sh\DTcomment{prepara a placa: FPGA, servidor, marca}.
.2 instala-quartus.sh\DTcomment{Quartus 20.1 para todas as contas de um PC}.
.2 deploy.sh, gera\_proveniencia.sh, PROVENIENCIA.sha256.
.2 Python/\DTcomment{o aplicativo}.
.3 morphe\_app.py\DTcomment{janela principal}.
.3 *\_window.py\DTcomment{uma janela por operação}.
.3 morphe\_protocol.py\DTcomment{cliente TCP}.
.3 dsp\_core.py, blocos.py\DTcomment{ponto fixo, referência, sinais longos}.
.3 fir\_design.py, iir\_design.py\DTcomment{projetistas de filtros}.
.3 ferramentas/\DTcomment{testes contra a placa, figuras, vetores}.
.2 C/\DTcomment{o servidor da placa}.
.3 morphe\_server.c, morphe\_protocol.h, morphe\_config.h.
.3 hps\_0.h\DTcomment{gerado do Platform Designer}.
.3 autostart/.
.2 Quartus/\DTcomment{o hardware}.
.3 ghrd\_top.v\DTcomment{topo: instancia os aceleradores}.
.3 conv1d.v, fft\_wrapper.v, iir\_*.v\DTcomment{os aceleradores}.
.3 soc\_system.qsys, soc\_system.qsf.
.3 tb\_iir\_*.v, roda\_tb\_iir.sh\DTcomment{testbenches}.
.3 output\_files/\DTcomment{o bitstream}.
.2 exemplos/\DTcomment{pacotes .mrph e sinais de exemplo}.
.2 docs/\DTcomment{documentação técnica e diário}.
.2 legado/\DTcomment{do projeto original, fora do produto}.
}
